% ============================================================
%  FINARA — Documentación Técnica de Módulos Principales
%  Trabajo Terminal — Ingeniería en Sistemas Computacionales
%  IPN UPIICSA
%  Fecha: Mayo 2026
%
%  Compilar con:
%    pdflatex documentacion_modulos.tex
%    pdflatex documentacion_modulos.tex   (segunda pasada para índice)
% ============================================================

\documentclass[12pt, letterpaper]{article}

% ── Codificación y lenguaje ──────────────────────────────────
\usepackage[utf8]{inputenc}
\usepackage[T1]{fontenc}
\usepackage[spanish, mexico]{babel}

% ── Geometría ────────────────────────────────────────────────
\usepackage{geometry}
\geometry{top=2.5cm, bottom=2.5cm, left=3cm, right=2.5cm}

% ── Tipografía y estilo ──────────────────────────────────────
\usepackage{lmodern}
\usepackage{microtype}
\usepackage{parskip}         % Espacio entre párrafos sin sangría
\setlength{\parindent}{0pt}

% ── Colores ───────────────────────────────────────────────────
\usepackage{xcolor}
\definecolor{azulipn}{RGB}{0, 66, 124}
\definecolor{verdefinara}{RGB}{19, 133, 61}
\definecolor{gristexto}{RGB}{60, 60, 60}
\definecolor{grisclaro}{RGB}{245, 247, 250}
\definecolor{bordesutil}{RGB}{200, 210, 220}
\definecolor{rojofinara}{RGB}{181, 24, 24}
\definecolor{codebg}{RGB}{248, 249, 250}
\definecolor{codeframe}{RGB}{210, 215, 220}
\definecolor{commentcolor}{RGB}{80, 120, 80}
\definecolor{keywordcolor}{RGB}{0, 80, 160}
\definecolor{stringcolor}{RGB}{140, 40, 120}

% ── Código fuente ─────────────────────────────────────────────
\usepackage{listings}
\lstdefinelanguage{Dart}{
  keywords={import, class, extends, implements, override, final, const,
            var, void, int, double, bool, String, List, Map, Future,
            async, await, return, if, else, for, while, null, true,
            false, new, this, super, get, set, required, late, static},
  keywordstyle=\color{keywordcolor}\bfseries,
  comment=[l]{//},
  morecomment=[s]{/*}{*/},
  commentstyle=\color{commentcolor}\itshape,
  stringstyle=\color{stringcolor},
  morestring=[b]',
  morestring=[b]",
  sensitive=true,
}
\lstdefinestyle{codestyle}{
  backgroundcolor=\color{codebg},
  basicstyle=\ttfamily\footnotesize,
  breaklines=true,
  breakatwhitespace=false,
  keepspaces=true,
  frame=single,
  rulecolor=\color{codeframe},
  numbers=left,
  numberstyle=\tiny\color{gray},
  numbersep=6pt,
  xleftmargin=1.2em,
  xrightmargin=0.5em,
  aboveskip=0.7em,
  belowskip=0.7em,
  showstringspaces=false,
  tabsize=2,
}
\lstset{style=codestyle, language=Dart}

% ── Tablas ────────────────────────────────────────────────────
\usepackage{booktabs}
\usepackage{array}
\usepackage{longtable}
\usepackage{tabularx}

% ── Gráficos y cajas ──────────────────────────────────────────
\usepackage{graphicx}
\usepackage{tcolorbox}
\tcbuselibrary{skins, breakable}

% Caja informativa (azul)
\newtcolorbox{infobox}[1][]{
  colback=azulipn!5, colframe=azulipn!50,
  fonttitle=\bfseries\small, title=#1,
  left=6pt, right=6pt, top=4pt, bottom=4pt,
  breakable,
}

% Caja de nota técnica (verde)
\newtcolorbox{notabox}[1][Nota técnica]{
  colback=verdefinara!5, colframe=verdefinara!50,
  fonttitle=\bfseries\small, title=#1,
  left=6pt, right=6pt, top=4pt, bottom=4pt,
  breakable,
}

% ── Listas ────────────────────────────────────────────────────
\usepackage{enumitem}
\setlist[itemize]{leftmargin=1.5em, topsep=3pt, itemsep=2pt}
\setlist[enumerate]{leftmargin=1.8em, topsep=3pt, itemsep=2pt}

% ── Encabezados y secciones ───────────────────────────────────
\usepackage{titlesec}
\usepackage{fancyhdr}

\titleformat{\section}
  {\Large\bfseries\color{azulipn}}{\thesection.}{0.6em}{}
  [\vspace{-0.3em}\textcolor{azulipn}{\rule{\textwidth}{0.5pt}}\vspace{0.2em}]

\titleformat{\subsection}
  {\normalsize\bfseries\color{gristexto}}{\thesubsection}{0.5em}{}

\titleformat{\subsubsection}
  {\small\bfseries\itshape\color{gristexto}}{}{0em}{}

\titlespacing*{\section}{0pt}{1.4em}{0.6em}
\titlespacing*{\subsection}{0pt}{1em}{0.4em}
\titlespacing*{\subsubsection}{0pt}{0.8em}{0.3em}

% ── Encabezado / pie ──────────────────────────────────────────
\pagestyle{fancy}
\fancyhf{}
\fancyhead[L]{\small\textcolor{azulipn}{\textbf{FINARA}} — Documentación Técnica}
\fancyhead[R]{\small\textcolor{gray}{IPN UPIICSA · 2026}}
\fancyfoot[C]{\small\thepage}
\renewcommand{\headrulewidth}{0.4pt}

% ── Hipervínculos ─────────────────────────────────────────────
\usepackage{hyperref}
\hypersetup{
  colorlinks=true,
  linkcolor=azulipn,
  urlcolor=azulipn,
  pdftitle={FINARA - Documentación Técnica},
  pdfauthor={Jesús Cadenas},
}

% ── Macros ────────────────────────────────────────────────────
\newcommand{\archivo}[1]{\texttt{\small\color{gristexto}#1}}
\newcommand{\endpoint}[1]{\texttt{\small\color{azulipn}#1}}
\newcommand{\campo}[1]{\texttt{#1}}


% ============================================================
\begin{document}

% ── PORTADA ──────────────────────────────────────────────────
\begin{titlepage}
  \centering
  \vspace*{1.5cm}

  % Logo institucional
  {\Large\bfseries\color{azulipn}
    INSTITUTO POLITÉCNICO NACIONAL\\[0.3em]
    \normalsize Unidad Profesional Interdisciplinaria de\\
    Ingeniería y Ciencias Sociales y Administrativas\\[0.2em]
    UPIICSA
  }\par
  \vspace{0.5cm}
  \textcolor{azulipn}{\rule{0.7\textwidth}{0.8pt}}\par
  \vspace{0.4cm}
  {\small Ingeniería en Sistemas Computacionales — Trabajo Terminal}\par

  \vspace{1.6cm}
  {\Huge\bfseries\color{verdefinara} FINARA}\par
  \vspace{0.3cm}
  {\large\color{gristexto} Aplicación Móvil de Finanzas Personales}\par
  \vspace{1.2cm}
  \textcolor{bordesutil}{\rule{0.5\textwidth}{0.4pt}}\par
  \vspace{0.8cm}

  {\LARGE\bfseries Documentación Técnica\\[0.2em]
   \large\normalfont Módulos Principales de la Aplicación Móvil}\par

  \vspace{1.4cm}
  \begin{tabular}{rl}
    \textbf{Integrante:}   & Jesús Cadenas \\[4pt]
    \textbf{Director:}     & Trabajo Terminal \\[4pt]
    \textbf{Framework:}    & Flutter SDK \textasciicircum 3.6.0 / Dart \\[4pt]
    \textbf{Backend:}      & Node.js + Express + MySQL \\[4pt]
    \textbf{OCR:}          & Tesseract + Google Cloud Vision + Claude \\[4pt]
    \textbf{Fecha:}        & Mayo 2026 \\
  \end{tabular}

  \vfill
  \textcolor{gray}{\footnotesize
    Este documento cubre únicamente los módulos más representativos
    de la aplicación. El sistema completo incluye módulos adicionales
    de autenticación, transacciones, notificaciones, perfil y onboarding.
  }
\end{titlepage}

% ── ÍNDICE ───────────────────────────────────────────────────
\tableofcontents
\newpage


% ============================================================
\section{Introducción General}
% ============================================================

\subsection{Objetivo de la aplicación}

FINARA es una aplicación móvil de finanzas personales desarrollada como
Trabajo Terminal en la carrera de Ingeniería en Sistemas Computacionales del
IPN UPIICSA. Su objetivo principal es proporcionar a los usuarios mexicanos
una herramienta digital integral para el registro, análisis y control de sus
finanzas personales, con énfasis en la automatización y la inteligencia
aplicada al gasto cotidiano.

\subsection{Problemática que resuelve}

La ausencia de hábitos financieros saludables en la población mexicana es
una problemática documentada: la mayoría de las personas no lleva un registro
de sus gastos, no establece presupuestos mensuales y desconoce con precisión
el estado de sus deudas. FINARA ataca este problema ofreciendo:

\begin{itemize}
  \item Registro automatizado de gastos mediante lectura OCR de tickets de compra.
  \item Control de presupuestos por categoría con alertas de exceso.
  \item Gestión financiera de deudas con cálculo de amortización francesa.
  \item Categorización inteligente de productos usando historial del usuario.
  \item Panel de inicio con resumen financiero visual del mes en curso.
\end{itemize}

\subsection{Enfoque y alcance del sistema}

La aplicación está orientada a usuarios con teléfonos Android que realizan
compras cotidianas en establecimientos que emiten tickets físicos. El flujo
principal consiste en:

\begin{enumerate}
  \item El usuario fotografía uno o varios tickets de compra.
  \item El sistema extrae texto mediante OCR con Tesseract y, como respaldo,
        Google Cloud Vision junto con Claude.
  \item Los productos detectados se categorizan automáticamente.
  \item El gasto queda registrado asociado a una categoría y fecha.
  \item Los presupuestos y alertas se actualizan en tiempo real.
\end{enumerate}

\begin{infobox}[Alcance de este documento]
  La aplicación FINARA cuenta con más módulos y funcionalidades, incluyendo
  autenticación, transacciones manuales, notificaciones push, perfil de
  usuario y proceso de onboarding. \textbf{Sin embargo, este documento
  documenta únicamente los cuatro módulos más relevantes y representativos
  del funcionamiento principal del sistema}: escaneo de tickets, categorías,
  presupuestos y deudas. Esto se realiza para mantener una documentación
  clara, resumida y enfocada en los componentes críticos de la aplicación.
\end{infobox}


% ============================================================
\section{Arquitectura General de la Aplicación}
% ============================================================

\subsection{Stack tecnológico}

\begin{longtable}{>{\bfseries\small}p{3.5cm} p{9.5cm}}
  \toprule
  Componente & Descripción \\
  \midrule
  \endhead
  Flutter SDK & Framework principal para la aplicación móvil multiplataforma.
    Versión \textasciicircum 3.6.0. Target: Android. \\[4pt]
  Dart & Lenguaje de programación tipado para el frontend móvil. \\[4pt]
  Node.js + Express & Backend REST API. Maneja autenticación JWT, lógica de
    negocio, invocación del pipeline OCR y persistencia. \\[4pt]
  MySQL & Base de datos relacional con tablas para usuarios, categorías,
    tickets, productos, presupuestos, deudas y alertas. \\[4pt]
  Python 3 (Tesseract) & Script \archivo{LecturaTickets\_v2.py} invocado
    como subproceso por Node.js. Ejecuta el pipeline OCR completo. \\[4pt]
  Google Cloud Vision & API de visión por computadora usada como respaldo
    cuando la confianza de Tesseract es insuficiente. \\[4pt]
  Claude Haiku & Modelo de lenguaje usado para estructurar el texto extraído
    por Vision en JSON de productos. \\[4pt]
  Firebase Crashlytics & Monitoreo de errores y excepciones en producción. \\[4pt]
  SharedPreferences & Almacenamiento local del token JWT y estado de sesión. \\
  \bottomrule
\end{longtable}

\subsection{Estructura del proyecto Flutter}

El proyecto sigue una arquitectura \textbf{feature-first}: cada módulo
funcional reside en su propio directorio dentro de \archivo{lib/features/},
con sus propias pantallas, servicios y widgets. El núcleo compartido se
ubica en \archivo{lib/core/}.

\begin{lstlisting}[language={}, numbers=none, frame=none, backgroundcolor=\color{white}]
lib/
├── core/
│   ├── constants/
│   │   ├── app_colors.dart       (paleta global de colores)
│   │   └── api_constants.dart    (URLs de endpoints)
│   ├── models/
│   │   └── finance_models.dart   (modelos: Budget, Debt, Category…)
│   └── services/
│       └── api_service.dart      (cliente HTTP centralizado)
└── features/
    ├── scan/          (escaneo OCR de tickets)
    ├── categories/    (gestión de categorías)
    ├── budgets/       (presupuestos mensuales)
    ├── debts/         (deudas e intereses)
    ├── transactions/  (gastos e ingresos)
    ├── home/          (pantalla principal)
    ├── auth/          (login, registro)
    ├── notifications/ (alertas y configuración)
    └── profile/       (perfil de usuario)
\end{lstlisting}

\subsection{Comunicación con el backend}

Toda la comunicación entre la aplicación móvil y el backend se realiza
mediante \textbf{API REST} sobre HTTPS. El cliente HTTP centralizado
\campo{ApiService} (singleton) gestiona:

\begin{itemize}
  \item Inyección automática del token JWT en el header \campo{Authorization}.
  \item Decodificación de respuestas en formato \campo{\{success, data\}}.
  \item Timeouts configurables: 30 segundos para peticiones normales,
        60 segundos para el endpoint OCR.
  \item Manejo de errores HTTP (400, 401, 403, 404, 409, 500).
  \item Subida de imágenes mediante \campo{multipart/form-data}.
\end{itemize}

La URL base del backend es: \endpoint{https://documents.digitaltax-analytics.com/finara/api}

\subsection{Manejo de estado}

La aplicación utiliza el patrón \textbf{StatefulWidget + setState} de Flutter
sin gestores de estado globales. Cada pantalla administra su propio estado
local. La capa de servicios (\campo{*Service}) actúa como intermediaria entre
la UI y el backend, abstrayendo la lógica HTTP.

\begin{notabox}
  No se utilizan paquetes como Provider, Riverpod o Bloc. El estado es local
  a cada widget. La comunicación entre pantallas se realiza mediante valores
  de retorno en \campo{Navigator.push} (\campo{result == true}) que disparan
  recargas en la pantalla padre.
\end{notabox}

\subsection{Paleta de colores del sistema}

\begin{tabular}{lll}
  \toprule
  \textbf{Nombre} & \textbf{Hex} & \textbf{Uso} \\
  \midrule
  \campo{prussianBlue} & \#0f172a & Color primario, textos, AppBar \\
  \campo{seaGreen}     & \#13853d & Ingresos, éxito, indicadores positivos \\
  \campo{mahoganyRed}  & \#b51818 & Gastos, errores, alertas críticas \\
  \campo{platinum}     & \#f1f5f9 & Fondo de pantallas \\
  \campo{white}        & \#ffffff & Superficies de tarjetas \\
  \bottomrule
\end{tabular}


% ============================================================
\section{Módulo de Escaneo de Tickets (OCR)}
% ============================================================

\subsection{Descripción general}

El módulo de escaneo es el componente más diferenciador de FINARA. Permite
al usuario fotografiar uno o varios tickets de compra y extraer automáticamente
los productos, precios, tienda y total mediante un pipeline de OCR de múltiples
etapas. El resultado alimenta directamente el registro de gastos y el control
de presupuestos.

\subsubsection{Flujo completo de operación}

\begin{enumerate}
  \item \textbf{Captura} (\campo{ScanTicketScreen}): El usuario activa la
    cámara trasera y captura entre 1 y 5 fotografías del ticket.
  \item \textbf{Previsualización} (\campo{ScanPreviewScreen}): Carrusel de
    imágenes con opción de eliminar, añadir más fotos o confirmar.
  \item \textbf{Procesamiento} (\campo{ScanProcessingScreen}): Las imágenes
    se suben al backend vía \campo{multipart/form-data}. El backend invoca el
    script Python OCR y devuelve el JSON estructurado.
  \item \textbf{Resultado} (\campo{ScanResultScreen}): El usuario revisa,
    edita o elimina los productos detectados antes de guardar el gasto.
\end{enumerate}

\subsection{Pantallas del módulo}

\subsubsection{ScanTicketScreen — Captura de imágenes}

Utiliza el plugin \campo{camera} para inicializar la cámara trasera con
resolución máxima (\campo{ResolutionPreset.max}) en formato JPEG.

\begin{lstlisting}
// Inicialización de cámara con permiso y resolución
Future<void> _initializeCamera() async {
  final cameras = await availableCameras();
  _controller = CameraController(
    cameras.first,
    ResolutionPreset.max,
    imageFormatGroup: ImageFormatGroup.jpeg,
  );
  await _controller.initialize();
}
\end{lstlisting}

Características relevantes:
\begin{itemize}
  \item Máximo de 5 imágenes por ticket para controlar el tamaño del payload.
  \item Visualización en tiempo real del encuadre con guía superpuesta
        (CustomPainter).
  \item Botón de flash (\campo{FlashMode.torch}) para ambientes oscuros.
  \item Galería inferior deslizable con las fotos capturadas y botón de
        eliminación individual.
\end{itemize}

\subsubsection{ScanPreviewScreen — Revisión previa}

Muestra las imágenes capturadas en un \campo{PageView} con navegación
manual. El usuario puede:

\begin{itemize}
  \item Eliminar la imagen actual con confirmación.
  \item Regresar a la cámara para añadir más fotos (respeta el límite de 5).
  \item Confirmar el lote de imágenes y pasar al procesamiento.
\end{itemize}

\subsubsection{ScanProcessingScreen — Ejecución OCR}

Es la pantalla de espera durante el procesamiento. Muestra 6 pasos animados
que orientan al usuario sobre el progreso:

\begin{enumerate}[label=\arabic*.]
  \item Preparando imágenes
  \item Subiendo al servidor
  \item Ejecutando OCR
  \item Extrayendo texto
  \item Identificando productos
  \item Calculando totales
\end{enumerate}

La petición al backend se realiza mediante \campo{TicketsService.scanTicket()}:

\begin{lstlisting}
// Subida multipart al endpoint POST /api/tickets/scan
Future<Map<String, dynamic>> scanTicket(List<String> imagePaths) async {
  return await _apiService.postMultipart(
    ApiConstants.ticketsScan,
    imagePaths,
  );
}
\end{lstlisting}

Si el OCR falla, la pantalla presenta un estado de error con botones para
reintentar o cancelar (regresa al inicio sin persistir nada).

\subsubsection{ScanResultScreen — Revisión y guardado}

Muestra los productos extraídos en una lista editable. El usuario puede:

\begin{itemize}
  \item Editar nombre, precio y categoría de cada producto.
  \item Eliminar productos no deseados.
  \item Añadir productos manualmente que no fueron detectados.
  \item Cambiar la fecha de la compra con un \campo{DatePicker}.
  \item Ver la confianza OCR del escaneo (indicador verde/naranja).
  \item Guardar el gasto completo.
\end{itemize}

\begin{lstlisting}
// Guardado del gasto con productos extraídos
Future<void> _saveExpense() async {
  await _expensesService.createWithProducts(
    idTicket:    _ticketId,
    nombreGasto: _expenseNameController.text,
    fecha:       _selectedDate,
    productos:   _extractedItems.map((item) => {
      'nombre_producto': item['name'],
      'precio':          item['price'],
      'id_categoria':    item['id_categoria'],
    }).toList(),
  );
}
\end{lstlisting}

\subsection{Pipeline OCR en el backend}

El script \archivo{ocr/LecturaTickets\_v2.py} implementa un pipeline de
6 etapas ejecutado como subproceso de Node.js:

\begin{enumerate}
  \item \textbf{Preprocesamiento} (OpenCV): upscaling a mínimo 1200 px,
        corrección de iluminación, deskew, binarización adaptativa gaussiana,
        denoising y sharpening.
  \item \textbf{OCR con Tesseract}: modo LSTM (\campo{oem 3}), bloque de
        texto (\campo{psm 6}), idiomas \campo{spa+eng}. Retorna texto y
        confianza promedio por palabra.
  \item \textbf{Evaluación de calidad}: dos compuertas independientes.
  \item \textbf{Respaldo Vision + Claude} (si aplica).
  \item \textbf{Parsing por regex}: extracción de productos, total y fecha.
  \item \textbf{Categorización}: pipeline de 3 pasos por producto.
\end{enumerate}

\subsubsection{Compuertas de calidad para activar el respaldo}

\begin{tabular}{llp{6cm}}
  \toprule
  \textbf{Compuerta} & \textbf{Condición} & \textbf{Descripción} \\
  \midrule
  1 — Técnica &
    \campo{conf < 0.72} O \campo{palabras < 15} &
    Baja confianza por palabra O texto escaso. Se activa antes del parsing. \\[4pt]
  2 — Estructural &
    Sin total Y sin productos &
    Tesseract pasó la Compuerta 1 pero el parsing no encontró datos útiles.
    Se activa después del parsing. \\
  \bottomrule
\end{tabular}

\subsubsection{Respaldo con Google Cloud Vision + Claude}

Cuando se activa cualquier compuerta, la función \campo{\_enhance\_extraction()}:

\begin{enumerate}
  \item Codifica la imagen en base64 y la envía a la API de Google Cloud Vision
        (\campo{DOCUMENT\_TEXT\_DETECTION}) con hint de idioma \campo{es+en}.
  \item Si se obtiene texto, lo envía a Claude Haiku con un prompt que
        estructura el ticket en JSON (\campo{store\_name, date, total, products}).
  \item El JSON enriquecido se deserializa al formato interno del pipeline.
  \item La confianza se establece en 0.95 si Vision fue exitoso.
\end{enumerate}

Si las claves API no están disponibles o cualquier paso falla, la función
retorna el texto original de Tesseract sin lanzar excepción.

\subsubsection{Categorización automática de productos}

Cada producto detectado pasa por un pipeline de 3 pasos:

\begin{enumerate}
  \item \textbf{categorias.json}: comparación del nombre del producto contra
        palabras clave organizadas por categoría. Match por substring exacto
        o similitud \campo{SequenceMatcher} $\geq 0.75$.
  \item \textbf{Historial del usuario}: si el Paso 1 no encontró categoría,
        se consultan los productos previos del mismo usuario en
        \campo{Productos\_Tickets} (filtrado por \campo{id\_usuario} vía join
        con \campo{Tickets} y validado en \campo{Categorias}). Match por
        substring o similitud $\geq 0.70$.
  \item \textbf{Fallback}: si ninguno encontró coincidencia, el producto se
        asigna a la categoría ``Otros'' del usuario con \campo{id\_categoria = null}.
\end{enumerate}

\begin{lstlisting}
# Pipeline de categorización — Paso 1 (categorias.json)
def _classify_by_json(name, categories):
    name_up = name.upper()
    for cat, keywords in categories.items():
        if cat == "OTROS": continue
        for kw in keywords:
            if kw.upper() in name_up:
                return cat          # Coincidencia exacta por substring
            if SequenceMatcher(None, name_up, kw.upper()).ratio() >= 0.75:
                return cat          # Coincidencia fuzzy
    return None                     # Sin coincidencia
\end{lstlisting}

\subsubsection{Métrica de confianza del ticket}

La función \campo{calculate\_confidence()} combina la confianza técnica
de Tesseract con señales estructurales del parsing:

\begin{longtable}{lrp{6.5cm}}
  \toprule
  \textbf{Componente} & \textbf{Peso} & \textbf{Justificación} \\
  \midrule
  \endhead
  \campo{ocr\_conf × 0.30}  & 30\% &
    Confianza promedio por palabra de Tesseract (normalizada a 0--1). \\[3pt]
  \campo{total\_found}       & +25\% / $-$15\% &
    El total es el campo crítico de todo ticket. \\[3pt]
  \campo{n\_products > 0}    & +20\% / $-$15\% &
    Los productos son el contenido principal. \\[3pt]
  \campo{store\_found}       & +10\% &
    Contextual; confirma que el documento es un ticket. \\[3pt]
  \campo{date\_found}        & +10\% / $-$5\% &
    Permite clasificación temporal del gasto. \\[3pt]
  \campo{word\_count < 10}   & $-$10\% &
    Imagen casi ilegible aunque las pocas palabras sean nítidas. \\[3pt]
  \campo{word\_count < 25}   & $-$5\% &
    Imagen con texto escaso (posible ticket parcial). \\
  \bottomrule
\end{longtable}

El valor final se normaliza al rango $[0.0, 1.0]$ con \campo{max(0, min(score, 1))}.

\subsection{Servicio de tickets}

\begin{tabular}{lp{8cm}}
  \toprule
  \textbf{Método} & \textbf{Descripción} \\
  \midrule
  \campo{scanTicket(imagePaths)} &
    POST multipart a \endpoint{/tickets/scan}. Timeout 60 s. \\
  \campo{getAll()} &
    GET \endpoint{/tickets}. Lista de tickets del usuario. \\
  \campo{getById(id)} &
    GET \endpoint{/tickets/:id}. Ticket con imagen en base64. \\
  \campo{getProducts(ticketId)} &
    GET \endpoint{/tickets/:id/products}. Productos con categoría. \\
  \campo{delete(id)} &
    DELETE \endpoint{/tickets/:id}. Elimina ticket y productos. \\
  \bottomrule
\end{tabular}

\subsection{Consideraciones técnicas}

\begin{itemize}
  \item \textbf{Idempotencia}: el backend genera un hash SHA-256 de las
        imágenes para evitar procesar el mismo ticket dos veces (error 409
        en duplicados).
  \item \textbf{Limpieza}: los archivos temporales de imagen se eliminan del
        servidor tras procesar el ticket.
  \item \textbf{Transacciones DB}: el ticket y sus productos se insertan en
        una transacción con \campo{rollback} ante cualquier falla.
  \item \textbf{Manejo de errores en Flutter}: si el OCR lanza excepción,
        la pantalla de procesamiento permite reintentar sin necesidad de
        fotografiar nuevamente el ticket.
\end{itemize}


% ============================================================
\section{Módulo de Categorías}
% ============================================================

\subsection{Descripción general}

El módulo de categorías provee la estructura de clasificación que sustenta
todos los demás módulos. Cada gasto, producto de ticket, presupuesto e ingreso
queda asociado a una categoría del usuario. El sistema distingue dos tipos:
\textbf{predefinidas} (creadas durante el onboarding) y \textbf{personalizadas}
(creadas libremente por el usuario).

Las categorías predefinidas por defecto son: Alimentos, Entretenimiento,
Transporte, Ahorro y Ropa.

\subsection{Modelo de datos}

\begin{lstlisting}
class Category {
  final int    id;
  final int    idUsuario;
  final String nombreCategoria;
  final String tipoCategoria;  // 'predefinida' | 'personalizada'
  final String? icono;         // Nombre del icono (ej. 'restaurant')
  final String? color;         // Color hex (ej. '#13853d')

  IconData getIconData();      // Mapea icono a FlutterIcon
  Color    getColor();         // Parsea hex a Color de Flutter
}
\end{lstlisting}

El método \campo{getColor()} parsea cadenas hexadecimales en formato
\campo{\#RRGGBB} o \campo{0xFFRRGGBB}. Si el formato no es reconocido,
retorna un color neutro por defecto (\campo{Colors.blueGrey}).

\subsection{Pantallas}

\subsubsection{CategoriesScreen — Listado}

Muestra las categorías en dos secciones diferenciadas:

\begin{itemize}
  \item \textbf{Predefinidas}: grid de 3 columnas con icono y nombre.
        No son editables ni eliminables.
  \item \textbf{Personalizadas}: lista con icono, nombre y botón de edición.
        El usuario puede crear, editar y eliminar las suyas.
\end{itemize}

\subsubsection{NewCategoryScreen — Creación}

Formulario con tres campos:

\begin{enumerate}
  \item Nombre de la categoría (obligatorio, máx. 50 caracteres).
  \item Icono: selector de cuadrícula con 20+ opciones predefinidas.
  \item Color: selector de cuadrícula con 12 colores disponibles.
\end{enumerate}

\subsubsection{EditCategoryScreen — Edición}

Igual que la creación, pero precarga los valores actuales de la categoría
recibida como argumento de navegación.

\subsection{Servicio de categorías}

\begin{tabular}{lp{8.5cm}}
  \toprule
  \textbf{Método} & \textbf{Descripción} \\
  \midrule
  \campo{getAll()} &
    GET \endpoint{/categories}. Todas las categorías del usuario. \\
  \campo{getById(id)} &
    GET \endpoint{/categories/:id}. \\
  \campo{create(\{nombre, icono, color\})} &
    POST \endpoint{/categories}. Crea categoría personalizada. \\
  \campo{update(\{id, nombre, icono, color\})} &
    PUT \endpoint{/categories/:id}. \\
  \campo{delete(id)} &
    DELETE \endpoint{/categories/:id}. \\
  \campo{getPredefined()} &
    Filtra localmente \campo{tipoCategoria == 'predefinida'}. \\
  \campo{getCustom()} &
    Filtra localmente \campo{tipoCategoria == 'personalizada'}. \\
  \bottomrule
\end{tabular}

\subsection{Flujo funcional}

\begin{enumerate}
  \item El usuario navega a la pantalla de Categorías desde el módulo de
        Presupuestos o desde el menú de configuración.
  \item Ve el listado separado en dos secciones.
  \item Para crear: presiona el botón ``+'' y completa el formulario.
  \item El servicio llama al endpoint POST y recarga la lista en caso de éxito.
  \item La nueva categoría queda disponible inmediatamente para asignar en
        presupuestos, gastos y productos de tickets.
\end{enumerate}

\subsection{Consideraciones técnicas}

\begin{itemize}
  \item Las categorías predefinidas no pueden eliminarse desde el frontend
        (el botón de edición no aparece en la sección predefinidas).
  \item En caso de error al cargar, se muestra un estado de error con botón
        de reintentar en lugar de pantalla vacía.
  \item El servicio retorna lista vacía (no lanza excepción) ante cualquier
        error HTTP, para no bloquear formularios que dependan de las categorías.
\end{itemize}


% ============================================================
\section{Módulo de Presupuestos}
% ============================================================

\subsection{Descripción general}

El módulo de presupuestos permite al usuario asignar un límite de gasto
mensual por categoría. El sistema calcula automáticamente cuánto se ha
gastado en cada categoría durante el mes en curso y emite alertas cuando
el presupuesto está cerca de agotarse o ha sido superado.

Cada presupuesto cubre el rango de fechas del mes calendario actual
(del día 1 al último día del mes), calculado automáticamente tanto en el
frontend como en el backend.

\subsection{Modelo de datos}

\begin{lstlisting}
class Budget {
  final int      id;
  final int      idCategoria;
  final double   montoAsignado;
  final double   montoUtilizado;   // Actualizado por el backend
  final String   periodo;          // 'mensual' (fijo)
  final DateTime fechaInicio;
  final DateTime fechaFin;
  final bool     activo;

  // Propiedades calculadas
  double get porcentajeUtilizado =>
      (montoUtilizado / montoAsignado * 100).clamp(0, double.infinity);

  double get montoDisponible =>
      (montoAsignado - montoUtilizado).clamp(0, double.infinity);
}
\end{lstlisting}

\subsection{Pantallas}

\subsubsection{BudgetsScreen — Panel de presupuestos}

La pantalla principal del módulo presenta:

\begin{itemize}
  \item \textbf{Tarjeta resumen global}: presupuesto total asignado,
        total gastado y porcentaje de uso del mes.
  \item \textbf{Barra de progreso global}: cambia a rojo cuando el uso
        supera el 90\%.
  \item \textbf{Tarjetas por presupuesto}: icono de categoría, nombre,
        barra de progreso individual, monto gastado vs. asignado,
        porcentaje y botón de edición.
  \item \textbf{Estado vacío}: con ícono y llamada a la acción para crear
        el primer presupuesto.
\end{itemize}

\begin{lstlisting}
// Cálculo de totales en la pantalla
double get _totalBudget =>
    _budgets.fold(0, (s, b) => s + b.montoAsignado);

double get _totalSpent =>
    _budgets.fold(0, (s, b) => s + b.montoUtilizado);

double get _percentageUsed =>
    _totalBudget > 0 ? (_totalSpent / _totalBudget * 100) : 0;
\end{lstlisting}

\subsubsection{NewBudgetScreen — Creación de presupuesto}

El formulario presenta un grid de categorías disponibles para seleccionar
(2 columnas, con resaltado visual de la seleccionada) y un campo de monto.

Validaciones implementadas:
\begin{enumerate}
  \item Monto requerido y mayor a cero.
  \item Categoría seleccionada obligatoria.
  \item Verificación de presupuesto duplicado: si la categoría ya tiene
        un presupuesto activo, se muestra un aviso y no se crea uno nuevo.
\end{enumerate}

\begin{lstlisting}
// Verificación de presupuesto duplicado antes de crear
final existingBudget = await _budgetsService.getByCategory(
    _selectedCategoryId!);
if (existingBudget != null) {
  ScaffoldMessenger.of(context).showSnackBar(
    const SnackBar(content: Text(
        'Esta categoría ya tiene un presupuesto asignado')));
  return;
}
\end{lstlisting}

\subsubsection{EditBudgetScreen — Edición}

Permite modificar el monto asignado, la categoría y el estado activo/inactivo
de un presupuesto existente. Recibe el \campo{budgetId} como argumento de
navegación y carga los datos actuales del presupuesto.

\subsection{Servicio de presupuestos}

\begin{tabular}{lp{8cm}}
  \toprule
  \textbf{Método} & \textbf{Descripción} \\
  \midrule
  \campo{getAll()} &
    GET \endpoint{/budgets}. Incluye \campo{porcentaje\_usado} precalculado. \\
  \campo{getByCategory(categoryId)} &
    GET \endpoint{/budgets/category/:id}. Para validar duplicados. \\
  \campo{create(data)} &
    POST \endpoint{/budgets}. Fechas auto-calculadas por backend. \\
  \campo{update(id, data)} &
    PUT \endpoint{/budgets/:id}. \\
  \campo{delete(id)} &
    DELETE \endpoint{/budgets/:id}. \\
  \campo{getActiveBudgets()} &
    Filtra localmente \campo{activo == true}. \\
  \campo{getBudgetProgress(budget)} &
    Calcula estado: safe / warning / danger / exceeded. \\
  \bottomrule
\end{tabular}

\subsection{Estados de progreso de presupuesto}

\begin{tabular}{llp{6cm}}
  \toprule
  \textbf{Estado} & \textbf{Umbral} & \textbf{Indicación visual} \\
  \midrule
  \campo{safe}     & $< 50\%$  & Barra azul \\
  \campo{warning}  & 50--80\%  & Barra naranja \\
  \campo{danger}   & 80--100\% & Barra roja \\
  \campo{exceeded} & $> 100\%$ & Barra roja + alerta \\
  \bottomrule
\end{tabular}

\subsection{Flujo funcional}

\begin{enumerate}
  \item El usuario accede al módulo de Presupuestos desde el
        \campo{BottomAppBar} en la pantalla principal.
  \item La pantalla carga todos los presupuestos activos del mes.
  \item El usuario puede ver el progreso individual de cada categoría.
  \item Si desea crear uno nuevo, selecciona la categoría y el monto límite.
  \item El backend auto-asigna las fechas del mes actual.
  \item Cuando un gasto nuevo se registra (OCR o manual), el backend
        actualiza el campo \campo{monto\_utilizado} del presupuesto
        correspondiente.
  \item Al rebasar el 80\% o el 100\%, el sistema de alertas genera una
        notificación para el usuario.
\end{enumerate}

\subsection{Consideraciones técnicas}

\begin{itemize}
  \item El campo \campo{montoUtilizado} es calculado y actualizado por el
        backend; el frontend solo lo lee.
  \item El mapa \campo{Map<int, Category>} en la pantalla permite búsquedas
        en O(1) sin recorrer la lista completa de categorías en cada render.
  \item Al regresar de la creación/edición con \campo{result == true},
        la pantalla recarga automáticamente la lista completa.
  \item El módulo llama a \campo{NotificationTriggers().checkAlertsOnce()}
        al cargar para verificar si hay alertas pendientes de presupuesto.
\end{itemize}


% ============================================================
\section{Módulo de Deudas}
% ============================================================

\subsection{Descripción general}

El módulo de deudas gestiona créditos e instrumentos de deuda con soporte
para la \textbf{fórmula de amortización francesa}. Permite registrar el
capital inicial, la tasa de interés anual, el plazo en meses y la fecha
del próximo pago. El backend calcula el pago mensual, la tasa mensual
equivalente y el saldo insoluto tras cada abono. El usuario puede realizar
pagos ordinarios y abonos extraordinarios a capital.

\subsection{Modelo de datos}

\begin{lstlisting}
class Debt {
  final int      id;
  final String   nombreDeuda;
  final double   capitalInicial;
  final double   tasaAnual;
  final double   tasaMensual;      // tasaAnual / 12
  final double   saldoInsoluto;    // Saldo pendiente post-pagos
  final double   pagoMensual;      // Calculado por backend
  final double   montoPagado;      // Acumulado de pagos
  final int      duracionMeses;
  final DateTime fechaSiguientePago;
  final bool     activa;

  // Propiedades calculadas en cliente
  double get porcentajePagado  => montoPagado / capitalInicial * 100;
  double get montoPendiente    => (capitalInicial - montoPagado)
                                    .clamp(0, double.infinity);
  double get interesTotal      => (pagoMensual * duracionMeses)
                                    - capitalInicial;
  bool   get tieneInteres      => tasaAnual > 0;
  int    get diasHastaPago     => fechaSiguientePago
                                    .difference(DateTime.now()).inDays;
}
\end{lstlisting}

\subsection{Fórmula de amortización francesa}

El pago mensual se calcula usando la fórmula estándar de cuota fija:

\[
  M = C \cdot \frac{r}{1 - (1 + r)^{-n}}
\]

donde $C$ es el capital inicial, $r$ es la tasa mensual ($r = \text{tasa\_anual} / 1200$)
y $n$ es el número de mensualidades. El frontend calcula esta fórmula en
tiempo real para mostrar una vista previa antes de guardar la deuda:

\begin{lstlisting}
// Cálculo del pago mensual con amortización francesa
double get _tasaMensual => _tasaAnual / 1200;

double get _pagoMensual {
  if (_capital <= 0 || _meses <= 0) return 0;
  if (_tasaMensual == 0) return _capital / _meses;
  final r = _tasaMensual;
  final n = _meses;
  return _capital * r / (1 - pow(1 + r, -n));
}
\end{lstlisting}

Para deudas sin interés (\campo{tasaAnual == 0}), el cálculo simplifica a
la división directa del capital entre el número de meses.

\subsection{Pantallas}

\subsubsection{DebtsScreen — Panel de deudas}

Muestra el estado consolidado de todas las deudas del usuario:

\begin{itemize}
  \item Tarjeta resumen: deuda total (naranja), total pagado (verde) y
        porcentaje pagado global.
  \item Por cada deuda: nombre, barra de progreso (rojo $<$ 50\%,
        naranja $\geq$ 50\%), monto pagado vs.\ total, días al próximo pago
        y botones de pago y edición.
  \item El botón de pago abre un diálogo modal donde el usuario ingresa
        el monto a pagar.
\end{itemize}

\subsubsection{NewDebtScreen — Creación}

El formulario captura los datos de la deuda y actualiza en tiempo real
una tarjeta de resumen financiero:

\begin{itemize}
  \item Nombre de la deuda.
  \item Capital inicial (monto del crédito).
  \item Plazo en meses (campo de texto libre, por defecto 12).
  \item Tasa de interés anual en porcentaje (por defecto 0\%).
  \item Fecha del próximo pago (DatePicker, por defecto +30 días).
\end{itemize}

La tarjeta de resumen en tiempo real muestra: tasa mensual equivalente,
pago mensual calculado e interés total del crédito.

\subsubsection{EditDebtScreen — Edición y abono a capital}

Además de editar los campos del formulario, incluye un botón exclusivo para
registrar \textbf{pagos anticipados a capital}. Este tipo de pago reduce el
saldo insoluto directamente sin afectar la cuota de intereses del período,
lo que disminuye las mensualidades restantes.

\subsection{Servicio de deudas}

\begin{tabular}{lp{8.5cm}}
  \toprule
  \textbf{Método} & \textbf{Descripción} \\
  \midrule
  \campo{getAll()} &
    GET \endpoint{/debts}. Todas las deudas del usuario. \\
  \campo{create(data)} &
    POST \endpoint{/debts}. El backend calcula \campo{pago\_mensual} y
    \campo{tasa\_mensual}. \\
  \campo{update(id, data)} &
    PUT \endpoint{/debts/:id}. \\
  \campo{registerPayment(id, monto)} &
    POST \endpoint{/debts/:id/pay}. Registra abono ordinario.
    Retorna deuda con \campo{saldo\_insoluto} actualizado. \\
  \campo{registerCapitalPayment(id, monto)} &
    POST \endpoint{/debts/:id/capital-payment}. Abono extraordinario
    a capital. \\
  \campo{delete(id)} &
    DELETE \endpoint{/debts/:id}. \\
  \campo{getActiveDebts()} &
    Filtra localmente \campo{activa == true}. \\
  \campo{getUpcomingPayments(\{days\})} &
    Filtra deudas con \campo{diasHastaPago $\leq$ days}. \\
  \campo{getDebtProgress(debt)} &
    Retorna estado: critical / warning / progress / almost / paid. \\
  \bottomrule
\end{tabular}

\subsection{Estados de progreso de deuda}

\begin{tabular}{llp{6cm}}
  \toprule
  \textbf{Estado} & \textbf{Umbral} & \textbf{Indicación visual} \\
  \midrule
  \campo{critical} & $< 25\%$ pagado    & Rojo \\
  \campo{warning}  & 25--50\% pagado    & Naranja \\
  \campo{progress} & 50--75\% pagado    & Amarillo \\
  \campo{almost}   & 75--100\% pagado   & Verde claro \\
  \campo{paid}     & $= 100\%$          & Verde \\
  \bottomrule
\end{tabular}

\subsection{Flujo funcional}

\begin{enumerate}
  \item El usuario accede al módulo de Deudas desde el \campo{BottomAppBar}.
  \item El panel carga y muestra todas las deudas activas con su estado.
  \item Para registrar un pago: botón verde en la tarjeta de deuda $\to$
        diálogo modal con campo de monto $\to$ llamada al endpoint
        \endpoint{/debts/:id/pay} $\to$ recarga de la lista.
  \item Para abono a capital: editar deuda $\to$ botón de abono a capital
        $\to$ diálogo con monto $\to$ llamada a
        \endpoint{/debts/:id/capital-payment}.
  \item El módulo verifica alertas de pagos próximos al cargar
        (\campo{NotificationTriggers}).
\end{enumerate}

\subsection{Consideraciones técnicas}

\begin{itemize}
  \item \textbf{Cálculo en frontend vs.\ backend}: la fórmula de amortización
        se ejecuta en tiempo real en el cliente para la vista previa, pero el
        backend recalcula y persiste los valores canónicos al guardar.
  \item \textbf{Validaciones}: el monto de pago debe ser mayor que cero y
        no puede exceder el saldo pendiente (\campo{montoPendiente}).
  \item \textbf{Deuda sin interés}: la fórmula degenera a división simple;
        la UI oculta la línea de ``Interés total'' cuando la tasa es cero.
  \item \textbf{Plazo variable}: el campo de plazo es un \campo{TextField}
        libre (no stepper) para facilitar el ingreso de valores grandes.
  \item Los días hasta el próximo pago se calculan en el cliente con
        \campo{DateTime.difference} para evitar tráfico innecesario.
\end{itemize}


% ============================================================
\section{Conclusiones}
% ============================================================

\subsection{Beneficios de la aplicación}

FINARA resuelve la necesidad de registro financiero cotidiano con una
propuesta que reduce la fricción del usuario al mínimo: fotografiar un
ticket es más rápido y menos propenso a errores que ingresar cada compra
manualmente. La combinación de OCR con Tesseract y el respaldo inteligente
con Google Cloud Vision y Claude garantiza una extracción de datos robusta
incluso en condiciones de imagen difíciles.

La categorización automática, reforzada con el historial de compras del
propio usuario, mejora con el uso: cuantos más tickets se escanean, más
precisa se vuelve la asignación de categorías para ese usuario.

\subsection{Automatización financiera}

Los tres pilares de automatización del sistema son:

\begin{enumerate}
  \item \textbf{Extracción OCR}: conversión automática de texto físico de
        un ticket en registros digitales estructurados.
  \item \textbf{Rollover mensual}: los presupuestos e ingresos recurrentes
        se renuevan automáticamente al inicio de cada mes sin intervención
        del usuario (job cron en el backend).
  \item \textbf{Alertas inteligentes}: el sistema monitorea en segundo plano
        el uso de presupuestos y la proximidad de pagos de deudas, enviando
        notificaciones proactivas.
\end{enumerate}

\subsection{Arquitectura modular y escalabilidad}

El sistema está diseñado bajo una \textbf{arquitectura modular}: cada módulo
funcional (escaneo, categorías, presupuestos, deudas) opera de forma
independiente con su propio conjunto de pantallas, servicios y endpoints.
Esto permite:

\begin{itemize}
  \item Agregar nuevos módulos (ej.\ inversiones, metas de ahorro) sin
        modificar el código existente.
  \item Reemplazar el motor OCR (Tesseract, Vision, otro) editando
        únicamente el script Python sin tocar el frontend.
  \item Migrar el estado local a un gestor global (Riverpod, Bloc) en el
        futuro sin reescribir la lógica de negocio.
  \item Escalar el backend de forma horizontal dado que cada request es
        stateless (autenticación por JWT).
\end{itemize}

\subsection{Posibles mejoras futuras}

\begin{itemize}
  \item \textbf{Soporte iOS}: el stack Flutter permite compilar para iOS con
        ajustes mínimos en los plugins de cámara y permisos.
  \item \textbf{Modo offline}: cachear presupuestos y gastos localmente con
        SQLite para operar sin conexión y sincronizar al reconectar.
  \item \textbf{Análisis predictivo}: usar el historial de gastos para
        proyectar el gasto del mes e identificar patrones de consumo.
  \item \textbf{Exportación de datos}: generación de reportes PDF/Excel
        con el resumen financiero mensual o anual.
  \item \textbf{Multi-moneda}: soporte para usuarios que manejan cuentas
        en diferentes divisas.
  \item \textbf{Mejora del OCR}: fine-tuning del modelo Tesseract con
        datasets de tickets mexicanos para aumentar la precisión de la
        confianza por palabra.
  \item \textbf{Integración bancaria}: conexión con APIs de open banking
        para importar movimientos automáticamente y reducir el registro
        manual.
\end{itemize}

\vspace{1.2cm}
\begin{center}
  \textcolor{bordesutil}{\rule{0.5\textwidth}{0.4pt}}\\[0.5em]
  {\small\color{gray}
    Documentación generada — Mayo 2026 \\
    FINARA · IPN UPIICSA · Trabajo Terminal
  }
\end{center}

\end{document}
